\chapter{Discussion}

\section{Interpretation of the Main Findings}
\label{sec:discussion_main_findings}

This thesis investigated complementary volumetric and slice-wise approaches to artifact detection in low-field fetal diffusion magnetic resonance imaging (dMRI). The experimental findings indicate that automated fetal-brain localization and quality classification can be connected in functional processing pipelines, but that successful operation of an individual model does not establish reliable quality assessment of an entire acquisition. The relevant distinctions are between segmentation overlap and anatomically useful crops, slice-level predictions and dynamic-level decisions, and retrospective computational feasibility and acquisition-time quality control.

The developed 3D framework combined nnU-Net v2 segmentation, brain-centered volume extraction, and 3D ResNet classification. Its recovered 600-dynamic prediction export showed 93.5\% accuracy and 100\% Artifact sensitivity. These figures demonstrate agreement with the recorded annotations, not independently validated generalization: the executed checkpoint, training exposure, and subject-level evaluation independence of the export could not be fully established (Section~\ref{sec:results_3d_classification}). The volumetric representation retains relationships across slices and directly corresponds to the original dynamic-level quality labels, but the apparent performance advantage cannot be assigned to spatial context without a controlled, independent comparison.

The 2D framework combined MobileNetV3-U-Net with EfficientNet-B0 to evaluate brain-centered individual slices. MobileNetV3-U-Net achieved a reported test Dice of 0.8620 on 3,168 slices from four held-out subjects, supporting its use as a fetal-brain localization component. The high typical overlap nevertheless concealed occasional empty or inaccurate masks, disproportionately observed among slices inherited from Artifact-labelled dynamics. Standalone EfficientNet-B0 classification was much less sensitive to inherited positive labels: 26.2\% on the complete oracle-crop test set and 23.5\% on the successful predicted-crop subset. Because even reference-mask-assisted cropping yielded limited sensitivity, localization failures cannot alone explain the difficulty of the slice-classification task. The dynamic-to-slice transfer of quality labels is a further important source of uncertainty.

An additional finding concerned the information used to reach dynamic-level decisions. On 88 test dynamics, the original 2D decision rule achieved 16.7\% Artifact sensitivity; a balanced failure-aware rule reached 66.7\% while retaining 94.3\% specificity. Most positive detections under the latter rule arose from unavailable predicted brain crops rather than classifier-score triggers (Section~\ref{sec:results_dynamic_decisions}). This suggests that segmentation outcomes can contribute to retrospective quality assessment, although crop failure is not itself proof of a visible artifact and the small validation population limits confidence in threshold stability.

The matched 71-dynamic comparison showed distinct error patterns: the recorded 3D predictions were more sensitive to Artifact-labelled acquisitions, whereas the 2D failure-aware decisions produced fewer false positives. The four-subject sample, unresolved 3D training exposure, and differences in the decision mechanisms prevent a definitive ranking of model generalization. Likewise, millisecond-scale offline 2D inference is promising for processing individual images, but does not establish scanner-integrated real-time detection, measured early stopping, or a controlled end-to-end speed advantage. Together, these findings establish the principal research contribution as a critical assessment of the capabilities, dependencies, and evidential boundaries of two implemented quality-assessment approaches rather than a validated clinical deployment.

\section{Three-Dimensional Artifact-Detection Pipeline}
\label{sec:discussion_3d}

\subsection{Strengths of Volumetric Analysis}
\label{subsec:discussion_3d_strengths}

The central architectural strength of the developed 3D pipeline is its representation of a complete diffusion dynamic. Three-dimensional convolutions can describe features extending through neighbouring image planes, whereas a classifier receiving an isolated 2D crop cannot directly evaluate such relationships. This distinction is potentially important where degradation manifests as spatially distributed signal changes or inconsistencies across the imaged brain. It is a plausible benefit of volumetric analysis, however, rather than a causal effect established by the present experiments.

The combination of nnU-Net v2 localization and brain-centered cropping also provided an anatomically focused route to quality classification. The segmentation component was designed to restrict the downstream ResNet input to the fetal-brain region while preserving original MRI intensities within the extracted bounding box. This modular structure distinguishes anatomical localization from quality assessment and allows each stage to be inspected separately. No controlled full-field-versus-brain-crop ablation was recovered, so the magnitude of any classification benefit from region-of-interest extraction remains unknown.

A further advantage is the correspondence between annotation and prediction unit. Original Good and Artifact labels were attached to complete diffusion dynamics, and the 3D ResNet likewise produced one decision per dynamic. Unlike the slice-wise classifier, it did not need to assign an identical dynamic label to individual images or aggregate independently estimated slice scores. This alignment avoids one particular weak-supervision mismatch, without eliminating uncertainty in the underlying expert quality annotations.

The recovered 3D classification export and matched evaluation provide evidence that the implemented volumetric classifier can produce discriminative dynamic-level decisions on the available records. In the 600-record export, all 92 Artifact-labelled dynamics were identified, alongside 39 false-positive decisions on Good-labelled dynamics. The latter errors demonstrate that a highly sensitive operating point can still impose a cost in unnecessary quality alerts. Interpretation of both outcomes must remain conditional on the unknown overlap with model development data; the export is not evidence of perfect detection on unseen fetuses.

\subsection{Limitations of the 3D Approach}
\label{subsec:discussion_3d_limitations}

The most consequential limitation of the 3D evidence is incomplete evaluation independence. The documented nnU-Net allocation contained subjects represented in both training and test entries, and the 3D ResNet implementation partitioned dynamic records without explicit grouping by subject. Repeated dynamics from one fetus share anatomical and acquisition-related characteristics; placing them across partitions can yield overly optimistic performance estimates. The executed classifier checkpoint and its relation to all 600 exported predictions were not fully recovered, so the recorded confusion matrix must not be interpreted as independently held-out performance.

The segmentation stage presents a separate limitation. Conventional volumetric cropping requires a usable foreground mask; an empty mask prevents ordinary crop construction, whereas a nonempty but incorrectly localized mask may produce a misleading classifier input. Aggregate, matched held-out nnU-Net segmentation metrics and a complete candidate-to-prediction processing log were unavailable. Accordingly, the number of exported ResNet predictions does not establish the success rate or accuracy of the full segmentation-to-classification pipeline. Failures observed with the earlier custom 3D U-Net cannot be used as quantitative failure estimates for the later nnU-Net model.

In the recorded 3D predictions, false-positive Artifact classifications affected Good-labelled dynamics, while no false negatives appeared in the recovered export. Such an error distribution could lead to unnecessary review or reacquisition if the same decision behaviour persisted in practice. Because matched original images, segmentation outputs, crops, and classification errors were not consistently available, the causes of these false positives cannot be reliably assigned to motion, anatomy, or localization error.

A final limitation concerns decision timing. The 3D pipeline requires a completed reconstructed volume before standard segmentation, crop extraction, and classification. Historical offline processing reports described approximately 15--25 seconds per scan and 17--20 seconds per volume, identifying volumetric segmentation as an important computational component. Their hardware and timing boundaries were not verified as equivalent, and a separate 5.901\,ms prepared-crop ResNet benchmark omitted segmentation and initial crop construction. These measurements therefore do not yield one reproducible end-to-end latency. The folder-watching and Streamlit demonstrations established post-file-availability automation, not analysis of a live scanner stream or automatic interruption and reacquisition. These constraints motivated investigation of a slice-wise alternative without implying that dimensionality alone resolves all pipeline failures.

\section{Two-Dimensional Artifact-Detection Pipeline}
\label{sec:discussion_2d}

\subsection{Strengths of Slice-Wise Processing}
\label{subsec:discussion_2d_strengths}

The 2D framework replaced the complete-volume prediction unit with a sequence of independently processed slices. MobileNetV3-U-Net performed fetal-brain localization, and EfficientNet-B0 classified a 224-by-224 brain-centered crop when a usable predicted mask was available. This reduced the image extent required for each inference operation and permitted the research to examine decisions based on selected slices rather than always waiting for a complete volumetric classifier input.

The segmentation network, with approximately 6.69 million parameters, reached substantial overlap on the subject-separated test population. Its recorded batch-size-one inference means of 2.748 and 2.759\,ms per slice demonstrated low model inference cost under the evaluated GPU configuration. A separate stage-specific report gave approximately 5.239\,ms for successful-slice preprocessing, segmentation, cropping, and classification. These measurements are not interchangeable and their minor numerical differences remain unresolved; nevertheless, they support the feasibility of computationally inexpensive offline slice processing.

The modular architecture also enabled an informative distinction between reference-mask-assisted (oracle) crops and automatically generated crops, and between a classifier decision and the absence of a classifier input. The 2D segmentation and classification partitions used disjoint recorded subject identifiers, reducing the specific leakage risk observed in the earlier 3D data allocation. The small number of independent test subjects still limits generalization claims, but the subject-based partition is a methodological strength of the later experimental design.

\subsection{Segmentation and Classification Challenges}
\label{subsec:discussion_2d_challenges}

Strong average segmentation overlap did not guarantee reliable localization in every slice. Across brain-present test slices, mean Dice was 0.9106 but median Dice was 0.9700, indicating a tail of poorer outcomes. Within the 1,529-slice classification test population, mean Dice was 0.9511 for slices from Good-labelled dynamics and 0.8081 for those from Artifact-labelled dynamics. The corresponding crop-failure counts were two of 1,258 and 28 of 271, respectively (Section~\ref{subsec:results_2d_segmentation_subgroups}). The concentration of failures among Artifact-labelled acquisitions is relevant, but those observations were clustered within a few subjects and do not identify visible artifact mechanisms in individual slices.

Two distinct failures matter downstream: an empty mask yields no classifier crop, while an anatomically inaccurate nonempty mask can still generate an input. The latter condition was observed in 21 of the 1,499 successful predicted crops with Dice below 0.5. Consequently, the reported 98.2\% overall predicted-crop availability does not mean that 98.2\% of crops were anatomically reliable (Section~\ref{sec:results_roi_reliability}). Dice can support retrospective analysis of such failures but cannot be used as a live acceptance criterion when reference masks are unavailable.

The EfficientNet-B0 results further show that crop reliability is not the entire problem. The complete oracle-crop test contained 200 false negatives against 271 inherited positive slice labels, despite using reference-mask-assisted inputs. The predicted-crop classifier also had limited sensitivity on its successful-input subset. Training loss decreased while validation loss increased in both recorded variants, a pattern consistent with limited generalization, although the precise contributions of overfitting, class imbalance, loss optimization, and label quality were not isolated. Since the two crop variants were trained separately and one excluded missing crops, differences in their metrics cannot be attributed solely to segmentation accuracy. Section~\ref{sec:discussion_segmentation_influence} examines that dependency specifically.

\subsection{Influence of Weak Slice-Level Labels}
\label{subsec:discussion_2d_weak_labels}

The 2D classification objective inherited Good or Artifact labels from complete diffusion dynamics. This made supervised slice training possible with the available annotations but did not establish that every selected slice from an Artifact-labelled dynamic visibly contained corruption. A model may therefore be penalized for classifying an apparently unaffected slice as Good, even if its prediction is anatomically plausible. Conversely, agreement with an inherited label does not independently demonstrate detection of a visible slice artifact.

This granularity mismatch limits the interpretation of the confusion matrices. An observed false negative against the inherited target cannot automatically be described as a missed visually confirmed artifact. Label noise is a credible contributor to the difficulty of training and evaluating EfficientNet-B0, but the present evidence does not quantify how many slice targets are incorrect or establish that weak supervision caused the observed sensitivity. Limited 2D spatial context and other model-development factors remain alternative explanations.

The class distribution further complicated model selection: positive inherited labels accounted for approximately 17.8\% of training slices, 6.7\% of validation slices, and 17.7\% of test slices. Such shifts affect prevalence-sensitive metrics and make a small positive validation set especially consequential for precision--recall-oriented checkpoint and threshold selection. A manually reviewed slice-level subset would allow visible corruption to be distinguished from an affected parent dynamic and would support a more direct evaluation of slice recognition. That additional annotation experiment was not completed in the documented study and is therefore future work rather than an achieved validation result.

\section{Comparison of Volumetric and Slice-Wise Approaches}
\label{sec:discussion_3d_2d_comparison}

\subsection{Spatial Context and Artifact Detection}
\label{subsec:discussion_comparison_context}

The 3D and 2D frameworks differ both in the anatomical information provided to their classifiers and in the way a dynamic-level decision is produced. The 3D ResNet operates on a brain-centered volume and predicts a quality class directly; the 2D EfficientNet-B0 operates on slices and requires subsequent aggregation or a failure-aware decision policy. The advantage of volumetric spatial context is therefore coupled in this study to differences in models, supervision, preprocessing, and decision rules.

The matched evaluation restricted the comparison to the same 71 dynamics from four subjects, with 59 Good and 12 Artifact annotations. This removes differences in evaluated dynamic identities, but not differences in model-training independence or the prediction mechanisms themselves. The 3D predictions identified all 12 positive dynamics; the 2D failure-aware configuration identified six. These observations are consistent with a possible benefit of volumetric context, but cannot be used to isolate it as the explanation (Section~\ref{sec:results_3d_vs_2d}).

The recorded correctness pattern also differed: both systems were correct on 52 dynamics, only 3D was correct on ten, only 2D was correct on nine, and neither was wrong on the same dynamic in that file. This shows nonidentical errors in the recovered predictions, not the performance of a tested hybrid or ensemble. In particular, it would be invalid to infer that combining the systems would provide error-free detection on independent acquisitions.

\subsection{Classification Reliability and Error Trade-offs}
\label{subsec:discussion_comparison_errors}

The matched accuracies were close: 87.3\% for 3D and 85.9\% for 2D, or 62 versus 61 correct dynamics. This one-observation difference in accuracy concealed a much larger divergence in error type. The 3D export yielded no false negatives and nine false positives; the 2D failure-aware decisions yielded six false negatives and four false positives. Corresponding Artifact sensitivities were 100.0\% and 50.0\%, while specificities were 84.7\% and 93.2\%, respectively.

These trade-offs matter because missed corrupted acquisitions may permit unsuitable data to proceed, whereas false alarms may lead to unnecessary review or reacquisition. The appropriate operating point depends on costs not measured in the present study. Accuracy alone is therefore inadequate for selecting a quality-control system under class imbalance. Moreover, the 2D result is for a composite failure-aware policy: five of its six true-positive decisions in the matched cohort were attributed to missing-crop triggers. It is not a head-to-head classifier-only comparison against the ResNet.

The matched cohort contained only 12 positive dynamics and four independent subjects, so subject clustering and small changes in error counts substantially affect estimated sensitivity. Neither subject-clustered uncertainty estimates nor formal paired significance tests were established. Most importantly, the evaluated 3D predictions belong to an export whose training exposure remains unresolved, and the precise frozen-rule provenance of the 2D matched decisions is not fully recoverable. The matched analysis is therefore exploratory and does not establish a reliable generalization ranking.

\subsection{Computational Requirements and Early Feedback}
\label{subsec:discussion_comparison_runtime}

A meaningful computational comparison must account for both the processing boundary and the point at which image data become available. The standard 3D pipeline requires a reconstructed complete dynamic and its segmentation-derived crop before ResNet classification. The 2D pipeline processes individual slices and can, in principle, accumulate evidence incrementally, although a dependable dynamic decision still requires an appropriate aggregation or stopping policy.

The available timings cover different quantities: historical complete-volume workflows of approximately 15--25 and 17--20 seconds, prepared-crop 3D classification of 5.901\,ms per volume, dedicated 2D segmentation inference of approximately 2.75\,ms per slice, and a reported successful-slice 2D processing time of 5.239\,ms. The prepared-volume benchmark omits nnU-Net and cropping, whereas the successful-slice 2D figure includes multiple processing stages; the historical experiments lack verified identical conditions. Consequently, no defensible end-to-end 3D-to-2D speedup factor follows from these values (Section~\ref{sec:results_computational_performance}).

The more relevant potential distinction is early access to a decision. Slice-wise processing need not wait for a completed volumetric classifier input. The recorded sequential demonstration implemented ordered inference and first-trigger logic on reconstructed NIfTI images, but did not establish the distribution of first detections, independent prospective sequential sensitivity, or measured acquisition-time savings. Computation alone also omits reconstruction, data transfer, scanner interface, and communication with an acquisition controller. Accordingly, low offline slice latency supports further engineering investigation, not a claim of demonstrated real-time scanner intervention. Speed and decision reliability must be evaluated jointly: an early but inaccurate alert is not an effective quality-control outcome.

\section{Influence of Brain Segmentation on Artifact Classification}
\label{sec:discussion_segmentation_influence}

The segmentation stage was not an independent endpoint of the proposed artifact-detection system. Its practical function was to identify a brain-centered region from which the downstream classifier could estimate image quality. This makes localization a dependency of classification: a good segmentation score is useful only insofar as the resulting crop retains the relevant anatomy and image characteristics, and an unavailable or misplaced crop can affect whether a slice is evaluated at all. The distinction between segmentation performance, crop availability, and classification outcome is therefore necessary when interpreting the 2D experiments.

The use of reference and predicted masks created two different pathways through this dependency. Reference-mask-assisted, or oracle, crops used the available corrected brain masks to define the region of interest. Predicted crops used MobileNetV3-U-Net outputs and thus included the practical consequences of automatic localization. The two pathways applied corresponding crop geometry, including the brain bounding box, a 15\% margin, square padding, and resizing. The margin was intended to preserve surrounding anatomical context and reduce the likelihood of excluding relevant brain tissue when the boundary estimate was imperfect. Nevertheless, using identical crop construction did not make the classifier inputs identical: differences in the mask position or extent could move, enlarge, or truncate the region supplied to EfficientNet-B0.

Two types of segmentation-dependent failure emerged. In 30 of the 1,529 selected test slices, empty predicted masks prevented crop generation entirely. These slices had no ordinary predicted-crop classifier output; they were not negative classifier predictions. A separate set of 21 slices produced nonempty predicted masks and classifier-ready crops despite segmentation Dice below 0.5. This second failure mode is potentially less visible to a pipeline that monitors only whether an input file was generated. The technical success of cropping does not establish that the expected fetal-brain anatomy was localized. The qualitative cases described in Section~\ref{subsec:results_crop_qualitative} illustrate both situations, but the original visual panel remains pending incorporation into the final thesis.

The oracle-versus-predicted analysis helps delimit what can be learned about downstream classification. Across all eligible test slices, the oracle pathway included 1,529 examples while the predicted pathway included only 1,499; therefore, their full-cohort metrics were not directly matched. Restricting the comparison to the 1,499 common slices removed that difference in observation identities. On this matched subset, oracle-crop classification achieved 79.3\% accuracy, 19.3\% Artifact sensitivity, and ROC-AUC 0.6440, while predicted-crop classification achieved 78.7\% accuracy, 23.5\% sensitivity, and ROC-AUC 0.6335 (Section~\ref{subsec:results_matched_crops}). The oracle variant had slightly higher accuracy and discrimination by ROC-AUC, whereas the predicted variant identified ten additional inherited positive slices while producing 20 additional false positives. Neither variant consistently dominated the other across metrics.

These results do not quantify the causal effect of segmentation errors on classification. Although the matched samples were the same, the oracle and predicted classifiers had been trained separately, and the recovered checkpoint-selection provenance was not complete for every variant. Differences could arise from training or decision behaviour as well as crop content. In addition, the matched evaluation excludes the 30 images with no predicted crop. It therefore examines classifier behaviour conditional on automatic localization succeeding, rather than evaluating all possible outcomes of the proposed automated pipeline. A strictly controlled estimate of localization impact would require evaluating a fixed classifier with systematically varied reference and predicted crops while explicitly accounting for failed localization.

The results also show why improving segmentation alone should not be assumed to solve artifact classification. On the full oracle-crop test set, where the classifier did not depend on the predicted segmentation masks, sensitivity to inherited Artifact slice labels was only 26.2\%. Thus, imperfect predicted localization cannot be the sole explanation for weak agreement with those targets. The uncertain relationship between an individual slice's appearance and its parent dynamic's quality label remains relevant, alongside the limitations of slice-wise spatial context and classifier development described in Section~\ref{subsec:discussion_2d_weak_labels}. Conversely, improving the slice classifier alone would not resolve the missing-input problem caused by empty masks.

An operational consequence is that segmentation failure and classification confidence need different interpretations. A missing crop is an observable failure of the localization-to-classification pathway, not a low artifact probability. The retrospective failure-aware dynamic decision rules used this distinction as an additional trigger, demonstrating that intermediate processing failures can influence overall quality decisions. Yet an empty mask might reflect image degradation, unfamiliar anatomy, low signal, or model limitations; the experiments do not establish which mechanism produced each failure. Likewise, low reference-mask Dice is informative during retrospective analysis but cannot serve as a prospective failure signal when no reference mask is available. Section~\ref{sec:discussion_dynamic_decisions} examines how available failure indicators were incorporated into dynamic-level decisions and the consequences of doing so.

The central implication is that segmentation should be assessed according to its downstream function, not solely according to pixel overlap. A reliable automated workflow must account for whether a crop exists, whether it plausibly contains the intended anatomy, and how uncertainty or failure is represented in the final quality decision. The present experiments document those dependencies and expose their limits, but do not establish a calibrated, independently validated anatomical crop-quality gate or the causal improvement in artifact detection that such a mechanism might provide.

\section{Dynamic-Level Decision Strategies and Failure Handling}
\label{sec:discussion_dynamic_decisions}

Dynamic-level decision-making is a distinct stage of the 2D framework: it combines EfficientNet-B0 slice probabilities and, where specified, information about unavailable predicted brain crops. Because the reference quality labels describe complete diffusion dynamics, the aggregation operator and operating rule determine whether otherwise identical slice outputs yield a Good or Artifact decision.

Mean, median, maximum, and recorded top-$k$ aggregators encode different assumptions about how artifact evidence is distributed through a dynamic. A mean may dilute a few high-scoring slices, a median is less sensitive to isolated extremes, and a maximum can respond to a single noisy prediction. Thus, these rules need validation at the dynamic level rather than selection from slice-level classification metrics alone.

The original dynamic rule identified only three of 18 Artifact-labelled dynamics in the 88-dynamic test set while correctly identifying 66 of 70 Good-labelled dynamics. Its sensitivity of 16.7\% despite 78.4\% accuracy illustrates that class prevalence can conceal clinically relevant false negatives; the complete operating-point comparison appears in Table~\ref{tab:results_failure_aware_policies}.

The validation search selected two failure-aware policies: median aggregation at threshold 0.93 (balanced) and the recorded mean of the five highest available scores at threshold 0.98 (recall-oriented). Both also flagged a dynamic if at least one predicted crop was unavailable. These thresholds apply to aggregated dynamic scores, not to the separate slice-wise stopping demonstration.

In the 88-dynamic retrospective test, the balanced rule increased Artifact sensitivity from 16.7\% to 66.7\% while retaining 94.3\% specificity. The gain corresponded to nine additional true positives without additional false positives, but this was a change in the composite decision rule rather than a retraining or improvement of EfficientNet-B0.

The trigger attribution explains the gain: ten of the balanced policy's twelve true-positive dynamic decisions arose from missing-crop flags, rather than classifier-score flags. Missing localization therefore supplied information associated with the inherited quality annotations, not proof that the classifier recognized a visible artifact. These records do not distinguish acquisition degradation from model limitations as the cause of each missing crop.

Missing-crop status can be observed automatically, unlike a Dice score that requires an unavailable reference mask during inference. However, the flag captures empty or unusable crops, not nonempty masks that localize the wrong anatomy; the latter can still yield misleading classifier inputs, as discussed in Section~\ref{sec:discussion_segmentation_influence}.

The recall-oriented rule detected the same twelve Artifact-labelled dynamics but increased false positives from four to ten, decreasing specificity from 94.3\% to 85.7\%. It therefore offered no positive-detection benefit in this test population. A threshold chosen to maximize validation recall is not guaranteed to maximize test recall.

The balanced rule was the strongest recorded operating point, not an independently established optimum. Only six positive dynamics were present in the 86-dynamic validation search, so each positive observation materially influenced threshold selection. Moreover, the recovered records do not establish that all subsequent evaluations were insulated from further inspection or manual adjustment. Threshold stability remains to be demonstrated on a larger independent cohort.

A missing crop must not be silently encoded as a low Artifact probability: that would confuse processing failure with evidence of acceptable image quality. Treating missingness as a separate outcome is a useful methodological choice, but any final automatic rejection policy would need independent validation of its false-alert consequences.

The sequential implementation extended the concept to ordered slices and configurable first-trigger stopping. Its individual-slice score threshold and failure-trigger option were distinct from the retrospective aggregation policies.

Crucially, the retrospective 88-dynamic results used available slice evidence from reconstructed data; they do not measure decisions reached using only images available up to an early stopping point. Verified first-trigger distributions, prospective partial-sequence accuracy, and acquisition-time savings were not recovered. Therefore, the observed balanced-rule gains cannot be interpreted as scanner-time early detection.

In summary, the study demonstrates that dynamic-level interpretation and processing-failure handling materially affect retrospective classification. For operational use, both the selected threshold and the time at which sufficient evidence becomes available must be validated independently.

\section{Methodological Limitations}
\label{sec:discussion_limitations}

The interpretation of the results is constrained by four cross-cutting features of the research design: limited independent subjects, uncertain annotation granularity, non-equivalent and partly reconstructed experiments, and offline-only computational evaluation. These concerns qualify the scope of the conclusions rather than undoing the recorded observations.

\subsection{Dataset Size and Generalization}

One of the principal limitations of the present study concerns the number of independent fetal subjects represented in the experimental populations. Although the available imaging data provided numerous diffusion dynamics and individual MRI slices, these observations were not statistically independent. Multiple dynamics acquired from the same fetus share anatomical characteristics, acquisition conditions, and potentially subject-specific imaging patterns. Likewise, neighbouring slices within a dynamic contain correlated anatomical information. Consequently, a large number of image observations does not necessarily correspond to an equally large number of independent biological samples.

The 2D segmentation and classification inventories comprised 15,600 and 9,113 slices, respectively, but only 24 recorded subject groups were partitioned as 16 training, four validation, and four test subjects. Even with disjoint recorded subject identifiers, performance on four held-out fetuses may be sensitive to the particular anatomy and acquisition characteristics represented. Slice-level sample size must not be confused with the number of independent subjects.

The available experiments did not establish subject-clustered confidence intervals or formal statistical uncertainty estimates for the principal model comparisons. Consequently, the reported performance differences should be interpreted as descriptive findings for the evaluated populations rather than precise estimates of expected performance across a broader fetal MRI population. The absence of such uncertainty estimates is especially important when comparing configurations that differ by only a small number of correct predictions.

Good-labelled observations predominated, with a different positive-class prevalence in the 2D training, validation, and test partitions. This affects precision, F1-score, average precision, and the stability of operating-point selection. Class-sensitive metrics improve transparency but do not compensate for scarce independent positive cases.

Another important consideration is the diversity of the imaging data. The developed methods were evaluated using the available low-field fetal dMRI acquisitions, and the documented imaging protocol was associated with a 0.55\,T MRI system. The experiments did not establish external validation across independently collected datasets, different scanner manufacturers, magnetic field strengths, acquisition protocols, or clinical institutions. Consequently, the models' robustness to such variations remains unknown.

No systematic subgroup validation was established by gestational age, fetal pathology, or acquisition settings. The results therefore cannot determine whether performance varies across those characteristics.

\subsection{Annotation Quality and Label Uncertainty}

Anatomical masks and quality labels provide different reference targets. The former describe fetal-brain location, whereas the latter classify complete diffusion dynamics. Their respective uncertainties must be considered separately.

The reference masks originated from automatic segmentation followed by manual correction, but independent inter-observer agreement was not established. Consequently, agreement with those references is not proof of absolute anatomical accuracy.

Background-only slices and the treatment of empty reference and predicted masks can influence aggregate Dice and IoU. The original per-mask evaluator outputs needed to verify all averaging conventions were not recovered, so those aggregate values remain reported rather than independently reconstructed.

The original inventory included an Uncertain category that was excluded from binary training and testing. This reduced ambiguity in the chosen binary targets but left the system unevaluated on uncertain acquisitions.

Dynamic-level Artifact labels were inherited by selected 2D slices without separate confirmation of visible corruption in each slice. The resulting label uncertainty can affect optimization and evaluation; its magnitude was not quantified. The mechanism and classification consequences are discussed in Section~\ref{subsec:discussion_2d_weak_labels}.

A further limitation is that the binary quality-label formulation does not characterize different artifact mechanisms separately. The available evaluation did not establish independently verified labels distinguishing motion-related signal loss, geometric distortion, or other acquisition-related degradation. As a result, the study could not determine whether particular artifact categories were more readily detected by the 3D or 2D models.

\subsection{Experimental Comparability and Evaluation Independence}

The experiments were developed sequentially, with different input units, architectures, preprocessing, training histories, and evaluation sets. They do not form one controlled comparison in which dimensionality is the only changing factor.

The early custom 3D U-Net lacks a verified independent segmentation test, while the four-record Random Forest experiment used incorrect subject-to-annotation matching. These outputs document exploratory development, not valid comparative generalization performance.

The later nnU-Net partition contained some of the same subjects across its train/test volume-entry allocations. For the 600-record 3D ResNet export, executed checkpoint and classifier training exposure remain unverified. Accordingly, the reported accuracy and sensitivity are export-level observations, not estimates of held-out-subject generalization.

The 2D partitions were subject-separated according to the recovered metadata, but the oracle and predicted classifiers were trained separately and their complete test populations differed by 30 unavailable predicted crops. Even the matched 1,499-slice analysis cannot isolate crop-source causality without holding the classifier fixed.

The interpretation of dynamic-level decision strategies was also limited by the available experimental provenance. Although the recovered implementation documented validation-based selection of aggregation methods and classification thresholds, the complete execution history did not conclusively establish that all operating rules were frozen before every subsequent evaluation. Repeated inspection of a test population during model or threshold development could produce optimistic estimates of generalization, even when the recorded final script separates validation and test operations.

The matched 71-dynamic comparison controlled observation identity, not training provenance, architecture, input eligibility, or failure-aware decision logic. The recorded differences in sensitivity and specificity must therefore not be attributed solely to volumetric versus slice-wise representation.

Run-specific package versions, executed checkpoints, full input-to-output logs, and the origin of some reconstructed prediction exports were incompletely retained. These gaps limit reproducibility even where saved metrics and confusion matrices are arithmetically consistent.

\subsection{Limitations of Computational and Real-Time Evaluation}

The timing studies were retrospective measurements on reconstructed MRI files, not prospective tests on an operating scanner. Their measurement boundaries differed across the reported 3D and 2D experiments.

The historical 3D workflow times included broader processing than the prepared-volume ResNet benchmark; the separate 2D reports measured segmentation alone or successful multi-stage slice processing. These values cannot support an end-to-end speedup factor without a common hardware configuration and identical timing boundaries.

The millisecond-scale successful-slice results exclude some processing failures and contain an unresolved discrepancy between summed stage means and reported totals. They are evidence of feasible offline computation, not a fully reconciled end-to-end latency distribution.

Hardware dependence represents another limitation. The dedicated 2D segmentation benchmark identified an NVIDIA RTX PRO 6000 Blackwell GPU partition, whereas complete hardware equivalence could not be established for all historical experiments. Inference latency may vary with hardware resources, memory availability, numerical precision, model-loading conditions, and input-processing implementations. The reported measurements therefore cannot be assumed to apply directly to other computational environments, including hardware deployed near an MRI scanner.

Scanner-to-decision latency would additionally include reconstruction, transfer, preprocessing, and communication with an acquisition controller. No live scanner integration or validated acquisition-time reaction loop was demonstrated.

The sequential NIfTI demonstration lacked verified first-trigger distributions, decision accuracy on partial acquisitions, and measured time saved through stopping. A software stopping condition is not evidence that the scanner can be stopped earlier.

No prospective evaluation established whether warnings would improve retained image quality, justify reacquisition, or reduce review effort. The costs of false-positive and false-negative quality decisions remain operationally unmeasured.

These limitations collectively restrict the scope of the study's conclusions: the recorded workflows and retrospective results are informative, but independent clinical generalization, complete reproducibility, and scanner-integrated decision performance remain unverified. Their resolution requires execution-linked experimental records, appropriately separated evaluation cohorts, and prospective acquisition-oriented testing.

\section{Implications for Automated Fetal MRI Quality Control}
\label{sec:discussion_implications}

The principal design implication is that fetal dMRI quality control should be treated as an integrated decision process rather than a single classifier output. The required operations include data availability, brain localization, assessment of whether the resulting crop is usable, artifact scoring, and a decision that makes processing uncertainty visible. The preceding sections establish the available evidence and its limits; here the emphasis is on requirements for future system design.

The representation should be selected with the intended decision time in mind. Volumetric analysis preserves relationships across the complete reconstructed dynamic; slice-wise analysis permits processing of individual images. A possible staged workflow could exploit both, but this study did not develop or validate such a combined system.

The failure-aware findings motivate an explicit distinction between three possible system states: acceptable image quality, suspected artifact, and insufficient evidence because processing or localization failed. This is a proposed design requirement, not a tested three-class classifier. In particular, a missing brain crop should not automatically be interpreted as verified corruption or silently recoded as Good.

The training target should correspond to the level at which a decision is required. Dynamic-level quality annotations are suitable for evaluating complete-acquisition decisions, while expert-reviewed slice-level artifact visibility would be necessary to validate the recognition of individual corrupted images. Collecting both could clarify when local deterioration is sufficient to classify an acquisition as unusable.

Operational decision rules should be frozen on validation data and assessed on truly independent subjects before deployment. Calibration, rejection of uncertain outputs, and false-alarm costs deserve explicit evaluation; no retrospective threshold in this thesis should be treated as a validated clinical operating point.

A false negative could allow degraded data to enter diffusion analysis, whereas a false positive could cause needless reacquisition. Choosing the sensitivity--specificity trade-off therefore requires clinical and acquisition-workflow input, rather than maximizing accuracy alone.

A sequential system must validate decisions using only information available at each potential stopping point. First-trigger timing, sensitivity before completing the dynamic, and false-trigger consequences would need to be measured prospectively, rather than inferred from retrospective aggregation or GPU inference times.

Practical deployment would also require a tested interface for reconstructed or intermediate scanner data, a predictable compute environment, communication of warnings, and an approved response to a quality alert. The present folder-based and NIfTI simulations did not evaluate these interfaces.

In summary, this work identifies concrete requirements for later acquisition-oriented quality control: reliable anatomical localization, evidence-aligned labels, decision rules tested on independent subjects, explicit uncertainty states, and measured scanner-to-decision performance. These requirements follow from the observed pipeline dependencies without implying that clinical deployment or real-time intervention has already been achieved.